\section{理论框架}

内禀软函数（$S_I$）可以从一个非单态轻赝标介子大动量形状因子的 QCD 因子化中得到，该介子由 $\pi=\overline{q}_2\gamma_5 q_1$ 组成，跃迁流由两个夸克双线性算符构成，其横向间隔固定为 $\vec{b}=(\vec{n}_\perp b_\perp, 0)$，
\begin{align}\label{eq:FF}
F(b_\perp,P^z)=\langle \pi(-\vec{P})|(\overline{q}_1\Gamma q_1)(\vec{b}) (\overline{q}_2\Gamma q_2)(0)|\pi(\vec P)\rangle_c.
\end{align}
这里 $q_{1,2}$ 是不同味的轻夸克场，$\vec{P}=(\vec{0}_{\perp},P^z)$。{为了提取软因子，算符与介子态的选取应使图~\ref{fig:formfactor} 中的四条线各为不同的味，如文献~\cite{Ji:2019sxk}所指出的。}最简单的情形对应于图~\ref{fig:formfactor} 中的收缩，它与所谓 connected insertion（连通插入）具有相同的拓扑。因此在式~(\ref{eq:FF}) 右端加上了下标 $c$。按此构造，disconnected insertion（非连通插入）在这种方案中不起作用，本文即采用该方案。

可以证明，式~(\ref{eq:FF}) 中定义的形状因子可因子化为准 TMDWF $\Phi$ 与内禀软函数 $S_I$~\cite{Ji:2019sxk,Ji:2020ect}
%
\begin{align}\label{eq:FF_factorization}
&F(b_\perp,P^z)= {S_I(b_\perp)}\\ &\times {\int_0^1 dx\, dx' H(x,x',P^z)\Phi^{\dagger}(x',b_\perp,-P^z)\, \Phi(x,b_\perp,P^z)}\nonumber
\end{align}
其中 $H$ 为{微扰}硬核。
%
准 TMDWF $ \Phi$ 是坐标空间关联函数的傅里叶变换，
\begin{align}
\label{eq:phi_def}&\phi(z,b_\perp,P^z)=\lim_{\ell\to\infty}\frac{\phi_\ell(z,b_\perp,P^z,\ell)}{\sqrt{Z_E(2\ell,b_\perp)}},\\
&\phi_\ell(z,b_\perp,P^z,\ell)\nonumber\\
&=\Big\langle 0\Big|\overline q_1\left(\frac{z }{2}n^z +\vec b\right)\Gamma_\Phi\,{\cal W}(\vec b,\ell)q_2\left(-\frac{z}{2}n^z \right)\Big| {\pi(\vec{P})}\Big\rangle\nonumber.
\end{align}
上式中 ${\cal W}(\vec b,\ell)$ 是类空钉形（staple-shaped）规范链，
\begin{align}
{\cal W}(\vec b,\ell)&= {\cal P}{\rm exp} \left[ ig_s\int_{-\ell}^{z/2} \textrm{d}s\ n^z\cdot  A(n^z s+b_\perp)\right] \nonumber\\
& \times {\cal P}{\rm exp} \left[ ig_s\int_{0}^{b_\perp } \textrm{d}s\ n_\perp\cdot  A(-\ell n^z +s  n_\perp)\right] \nonumber\\
& \times {\cal P}{\rm exp} \left[ ig_s\int_{-z/2}^{-\ell} \textrm{d}s\ n^z\cdot  A(n^z s)\right],
\end{align}
$n^z$ 与 $n_\perp$ 分别是 $z$ 方向与横向方向的单位矢量。
$Z_E(2\ell,b_\perp)$ 是尺寸为 $2\ell \times b_\perp$ 的类空矩形 Wilson 圈的真空期望值，它去除准 TMDWF 中的 pinch-pole 奇异性与 Wilson 线自能~\cite{Ji:2019sxk}。

由于内禀软函数的 UV 发散是相乘性的~\cite{Ji:2020ect}，格点上可计算的比值 $S_I(b_{\perp},1/a)/S_I(b_{\perp,0},1/a)$ 与 UV 重正化方案无关，其中参考距离 $b_{\perp,0}$ 取得足够小以便微扰计算。于是我们可通过
\begin{align}\label{eq:S_ratio}
S_{I,\overline{\rm MS}}(b_\perp,\mu)= \left(\frac{S_I(b_{\perp},1/a)}{S_I(b_{\perp,0},1/a)}\right)S_{I,\overline{\rm MS}}(b_{\perp,0},\mu)
\end{align}
得到 $\overline{\rm MS}$ 方案下的结果，其中 $S_{I,\overline{\rm MS}}(b_{\perp,0},\mu)$ 可微扰计算，例如
\begin{align}\label{eq:S_ratio_1loop}
S_{I,\overline{\rm MS}}(b_{{\perp}},\mu)=1-\frac{\alpha_sC_F}{\pi}\ln\frac{\mu^2 b_{{\perp}}^2}{4 e^{-2\gamma_E}}+{\cal O}(\alpha_s).
\end{align}

在本探索性研究中，我们只考虑式~(\ref{eq:FF_factorization}) 中的领头阶匹配，此时微扰核为 $H(x,x',P^z)=1/({2N_c})+{\cal O}(\alpha_s)$，与 $x$ 和 $x'$ 无关。
利用宇称变换下的 $\phi(0,b_\perp,-P^z)=\phi(0,b_\perp,P^z)$，我们得到
\begin{align}\label{eq:S_I_z=0}
S_I(b_\perp)=\frac{2N_c F(b_\perp,P^z)}{|\phi(0,b_\perp,P^z)|^2}+{\cal O}(\alpha_s, (1/P^z)^2),
\end{align}
其中忽略了有限 $P^z$ 带来的幂次修正。由于 $P^z$ 与介子的快度相关联，下面我们用增强因子 $\gamma\equiv E_{\pi}/m_{\pi}$ 代替它。式~(\ref{eq:S_ratio}) 可写为
\begin{align}\label{eq:S_ratio_z=0}
S_{I,\overline{\rm MS}}(b_\perp,\mu)&=\frac{F(b_\perp,P^z)}{F(b_{\perp,0},P^z)}\frac{|\phi(0,b_{\perp,0},P^z)|^2}{|\phi(0,b_\perp,P^z)|^2}\nonumber\\
&+{\cal O}(\alpha_s, \gamma^{-2})\, .
\end{align}
%
由于只保留领头阶贡献，上式右端的比值与重正化标度 $\mu$ 无关。

%To our knowledge, the quasi-TMDWF has not been studied before on the lattice. Once it is known, it
{另一方面，准 TMDWF} 可用来提取 Collins-Soper 核 $K$，方法类似于文献~\cite{Ebert:2018gzl}
\begin{align}
&K(b_\perp,\mu)=\frac{1}{\ln(P_1^z/P_2^z)}\ln\left|\frac{C(xP_2^z,\mu) \Phi_{\overline{\rm MS}}(x,b_\perp,P_1^z,\mu)}{C(xP_1^z,\mu) \Phi_{\overline{\rm MS}}(x,b_\perp,P_2^z,\mu)}\right|\label{eq:CS_kernel}\\
&\quad=\frac{1}{\ln(P_1^z/P_2^z)}\ln\left|\frac{\int_0^1 \textrm{d}x \Phi(x,b_\perp,P_1^z)}{\int_0^1 \textrm{d}x \Phi(x,b_\perp,P_2^z)}\right|+{\cal O}(\alpha_s, \gamma^{-2})\nonumber\\
&\quad=\frac{1}{\ln(P_1^z/P_2^z)}\ln\left|\frac{\phi(0,b_\perp,P_1^z)}{\phi(0,b_\perp,P_2^z)}\right|+{\cal O}(\alpha_s, \gamma^{-2})\label{eq:CS_kernel_z=0}.
\end{align}
%
{在第二步中，同样只使用了领头阶匹配核 $C(xP^z,\mu)=1+{\cal O}(\alpha_s)$。}
$\Phi$ 的重正化因子被抵消。
%
在此近似下，由于只保留了领头项，与快度方案无关的 CS 核 $K$ 不依赖于 $\mu$。

式~(\ref{eq:S_ratio}) 与 (\ref{eq:CS_kernel}) 是严格的，可用于将来的精密研究；而式~(\ref{eq:S_ratio_z=0}) 与 (\ref{eq:CS_kernel_z=0}) 是这项开创性工作所用的领头阶近似。

\begin{table}[htbp]
\begin{center}
\caption{\label{table:cls} 数值模拟所用参数。第一行给出 2+1 味 clover 费米子 CLS 系综（编号 A654）的参数，第二行给出本计算所用的 A654 组态数目与价夸克π介子质量。}
\begin{tabular}{ccccccc}
\hline
 $\beta$ & $L^3\times T$  &a (fm)  & $c_{sw}$  & $\kappa^{\rm sea}_l$ & $m^{\rm sea}_{\pi} $(MeV)\\
  3.34 & $24^3 \times  48$ & 0.098 &  2.06686  & 0.13675 & 333 \\
\hline
&&& $N_{cfg}$  &  $\kappa^{v}_l$  & $m^v_{\pi} $ (MeV)  \\
  &&&  864  & 0.13622   & 547 \\
 \hline
\end{tabular}
\end{center}
\end{table}
